% SP2 solution v3.
% Key fix: "opposite signs exist" now uses the counting identity (SP1 item 7) instead
% of the broken surjectivity argument.  The rest is the same as v2.

Let \(P:=\operatorname{conv}(S)\). Since \(S\) is in strict convex position,
\(V(P)=S\) and \(|S|=n\). In the equality case we have from Subproblem~1:
\[
\operatorname{supp}(F)=\{\ell_s:s\in S\},\qquad F(\ell_s)=\sigma(s):=\varepsilon(t_s)\in\{\pm1\}.
\]

\medskip
\textbf{1. Opposite signs occur.}

By the counting identity (Subproblem~1, item~7),
\[
n_+ - n_- = n(k_1-k_2),
\]
where \(n_\pm = |\{s\in S:\sigma(s)=\pm1\}|\) and \(k_\pm = |\{t\in T:\varepsilon(t)=\pm1\}|\).

Since \(k_1\ge1\) and \(k_2\ge1\) (both signs present in \(T\)), we have \(k_1-k_2\in\{-1,0,1\}\).

\begin{itemize}
\item \textbf{Case~0} (\(k_1=k_2\)): \(n_+=n_-=n/2\).  Since \(n\ge3\) and \(n\) must be even,
  we have \(n\ge4\), \(n_+\ge2\), and \(n_-\ge2\).  \emph{In particular, both signs appear on \(S\).}
\item \textbf{Cases~\(\pm n\)} (\(k_1-k_2=\pm1\)): either \(n_+=n\) (all signs \(+1\)) or
  \(n_-=n\) (all signs \(-1\)).
\end{itemize}

\emph{For Subproblems~2--3 we assume Case~0.}  (Cases~\(\pm n\) are treated directly by
the algebraic arguments of Subproblem~4 without needing the chain decomposition:
the trinomial lemma handles \(n=3\); the Newton-polytope factorization handles \(n=4\).)

In Case~0 there exist \(s_a,s_b\in S\) with \(\sigma(s_a)=+1\) and \(\sigma(s_b)=-1\). \qed

\medskip
\textbf{2. Boundary-chain decomposition.}

Choose the pair \(s_a,s_b\) just found.  Since all inequalities defining ``\(s_a\) leftmost
and \(s_b\) rightmost'' are strict and open, we may additionally perturb the coordinate
direction so that no two vertices of \(P\) share the same \(x\)-coordinate and no edge of
\(P\) is vertical, while keeping \(s_a\) leftmost and \(s_b\) rightmost.

Removing \(s_a\) and \(s_b\) from \(\partial P\) splits the boundary into two arcs.
Their closures are the \textbf{upper chain}
\[
U=(v_1=s_a,\,v_2,\,\dots,\,v_{m_U}=s_b)
\]
and the \textbf{lower chain}
\[
L=(v_1=s_a,\,v_2,\,\dots,\,v_{m_L}=s_b),
\]
each a strictly \(x\)-monotone convex polygonal arc from \(s_a\) to \(s_b\).

\emph{Strict \(x\)-monotonicity:} For each \(c\in(x(s_a),x(s_b))\) the vertical line
\(\lambda_c\) meets \(P\cap\partial P\) in exactly two points (one on \(U\), one on \(L\)),
since \(P\) is convex and has no vertical edges.  Hence each chain meets \(\lambda_c\) in
exactly one point, so the \(x\)-coordinate is strictly increasing along each chain.

The vertex sets satisfy
\[
U\cup L=S,\qquad U\cap L=\{s_a,s_b\},
\qquad |U|+|L|=n+2.
\]
This proves~(2). \qed

\medskip
\textbf{3. Each chain contains both signs.}

Both \(U\) and \(L\) share the endpoints \(s_a\) (sign \(+1\)) and \(s_b\) (sign \(-1\)).
Hence each chain contains both signs. \qed
